\section{Main Results}

\begin{theorem}[Main Conclusion]
  \label{thm:finding-bimat-game-12_del-wsne-is-poly-time}
  For any constant \(\delta > 0\), there is a polynomial-time algorithm that finds a \((1/2 + \delta)\)-WSNE in any \(n \times n\) normalized bimatrix game \((\mathcal{R}, \mathcal{C})\) \cite[Theorem~1]{12WSNE}.
\end{theorem}

Note that the exponent of the ``polynomial time'' depends on \(\delta\); when \(\delta\) is a constant, the running time is polynomial in \(n\).

For games in which the payoff matrix is not given but only a payoff-query oracle is available, the paper also gives a query version:

\begin{theorem}[{Query-Efficient Version \cite[Theorem~2]{12WSNE}}]
  \label{thm:query-efficient-version}
  For any constants \(\varepsilon, \delta > 0\), with \(O\left( \frac{n \log n}{\varepsilon^4} \right)\) payoff queries, one can find a \((1/2 + 3\varepsilon + \delta)\)-WSNE in polynomial time with probability at least \(p_n^2\), where
  \[
    p_n = \left(1 - n^{-1 / 8}\right)\left(1 - \frac{2}{n}\right)^2.
  \]
\end{theorem}

The 2023 publication's factor \(p_n^2\) accounts for both calls to the zero-sum routine. When the subsequent small-support profile is obtained by random sampling with failure probability at most \(\eta\), a fully explicit implementation succeeds with probability at least \(p_n^2(1-\eta)\); for \(\eta=1/n\), this is \(p_n^2(1-1/n)\). This refinement and its unchanged complexity consequences are derived in Section~\ref{sec:source-corrections}.

In general, to compute a \((1/2 + \delta')\)-WSNE, one can take \(\varepsilon = \delta = \frac{\delta'}{4}\). Similarly, the polynomial time here also depends on \(\varepsilon\) and \(\delta\).
